%% kapa: exam=1 round=19 year=2008 period=1 form=A subject=civil_law subject_ko=민법 seq=18 number=18
%% source: Q-Net 감정평가사 기출문제 (cst00309) artlSeq=1001962 file=extracted/1008/1001962/2008년 제19회 감정평가사 1차시험 기출문제/1교시.hwp
%% conversion: kordoc 4.24.51 -> tools/gen1.py (auto); review=proofread-sonnet; flags=none
\begin{question}{18}
甲의 무권대리인 乙이 甲소유의 토지를 丙에게 매도하는 계약을 체결하였는데, 丙은 乙이 무권대리인임을 알지 못하였고, 乙의 대리행위가 표현대리에 해당되지 않는다. 다음 중 틀린 것은?
\begin{choices}
  \item 甲은 丙에 대하여 乙의 무권대리행위를 추인할 수 있다.
  \item 甲은 丙에 대하여 乙의 무권대리행위의 추인을 거절할 수 있다.
  \item 乙은 甲에 대하여 자신의 무권대리행위에 대한 추인여부의 확답을 최고할 수 있다.
  \item 추인이 있기 전에 丙은 甲에 대하여 계약을 철회할 수 있다.
  \item 추인이 있기 전에 丙은 乙에 대하여 계약을 철회할 수 있다.
\end{choices}
\end{question}
\answer{3}
